\section{Literature and previous research}
\label{sec:literature}

The introduction has shown that portfolio diversification depends on the stability of the
relationships between assets. However, empirical research indicates that correlations are
not constant over time and may change substantially during periods of financial stress.
An asset that appears to provide diversification under normal conditions may therefore
behave differently when markets experience high volatility and falling prices.

\citet{longin2001extreme} show that correlations between international equity markets
increase particularly during severe downturns. \citet{angchen2002asymmetric} document a similar
asymmetry for U.S. equity portfolios, where correlations are stronger after negative market
movements than after positive movements of comparable size. These findings imply that
unconditional correlations can understate the comovement investors face in adverse states.
For a diversification analysis, it is therefore not enough to ask whether assets
appear diversified on average. The relevant question is whether those relationships remain stable
when markets are under stress.

A defensible answer first requires a systematic definition of market stress. Historical crisis
dates can be selected manually, but such choices are partly subjective and may bias the
results. Markov-switching models offer a statistical alternative by allowing returns to move
between latent states with different means and variances \citep{hamilton1989new}. This
approach is well suited to financial markets because stress is usually characterized by high
volatility, lower returns, and persistent regime dynamics. Prior work shows that such regimes
matter for asset allocation because return distributions and dependence structures can differ
substantially across states \citep{ang2002international,guidolin2007asset}.

However, higher stress-period correlations must be interpreted carefully. \citet{forbesrigobon2002}
show that correlations can increase mechanically when volatility rises, even if the underlying
relationship between markets has not changed. They distinguish between contagion, meaning a
genuine strengthening of linkages after a shock, and interdependence, meaning strong existing
relationships observed in a more volatile environment. This distinction is central for our
paper. Raw correlations describe the comovement that investors actually experience, while
Forbes--Rigobon-adjusted correlations ask whether the increase remains after correcting for
heteroskedasticity. A stress correlation increase can therefore signal a realized loss of
diversification without necessarily proving structural contagion.

Pairwise correlations also provide only a partial view of diversification. Even if some pairs
remain stable, the portfolio as a whole may become more exposed to a common risk factor.
Principal component analysis captures this broader concentration by measuring how much of
total variation is explained by the leading components \citep{jolliffe2002pca}. A higher first
principal-component share in stress means that returns are increasingly driven by one common
direction, leaving fewer independent sources of risk. This idea is closely related to the
systemic-risk literature: \citet{kritzman2011principal} interpret a rising concentration of
variance in a few components as an absorption ratio, and \citet{billio2012econometric} show that
financial-system connectedness tightens during crises.

This study combines these strands of research in one empirical framework. It compares
regime-conditional correlations for a cross-asset universe and an international equity universe,
tests their sampling uncertainty, measures portfolio-level concentration with PCA, and applies the Forbes--Rigobon
adjustment to distinguish volatility-driven comovement from stronger evidence of contagion.